\documentclass[10pt,a4paper]{amsart}
\usepackage[margin=19mm]{geometry}
\usepackage{amsmath,amssymb,mathtools}
\usepackage[scr=rsfs,cal=euler]{mathalfa}
\usepackage{xfrac}
\usepackage[unicode,hidelinks,bookmarks=false]{hyperref}
\newcommand{\ie}{i.e.\ }
\newcommand{\cf}{cf.\ }
\newcommand{\resp}{respectively\ }
\newcommand{\Subsection}{\subsection}
\providecommand{\nobreakdash}{\nobreak\hbox{-}}
\setlength{\emergencystretch}{2em}
\multlinegap0pt
\usepackage{amssymb}
\usepackage{float}
\newtheorem{corollary}{Corollary}
\newcommand{\SN}{\mathrm{SN}}
\title{Right invariant Poisson Nijenhuis structures on Lie groupoids Correspondence and Classification}
\author{Ghorbanali Haghighatdoost}
\date{Source: arXiv:2411.17179v2 (CC BY 4.0)}
\begin{document}
\maketitle

\noindent\emph{Source and licence.} Right invariant Poisson Nijenhuis structures on Lie groupoids Correspondence and Classification. Version arXiv:2411.17179v2 (CC BY 4.0), distributed by arXiv under the Creative Commons Attribution 4.0 International licence, http://creativecommons.org/licenses/by/4.0/, as recorded in the arXiv licence metadata for this exact version and shown on the abstract page https://arxiv.org/abs/2411.17179v2. Full licence notice: \texttt{licenses/arxiv-2411.17179-CC-BY-4.0.txt}.\par\smallskip

\noindent\textit{Editorial scope.} Counted unit: the article's corollary cor:classification (source0.tex lines 465--497). The article's complete original proof of it is delivered with it (source0.tex lines 499--501, one \texttt{proof} environment of 3 lines). No mathematics is written, repaired or re-derived: the statement and the proof are the article's own lines, in the article's own notation, and no retained line is substituted or re-flowed.  No further source-local context is needed: every cross-reference of every retained line resolves inside the two delivered spans.  Nothing was omitted from any retained span, so the package carries no omission record.  No retained line contains a citation, so the wrapper carries no bibliography and no bibliography entry.  No retained line contains a figure environment, an image-inclusion command or drawing code, so no image is delivered and the package declares no asset dependency.  Editorial formatting only, in addition: an AMS \texttt{amsart} preamble that restates the article's own declarations and macros used by the retained lines, namely the environments \texttt{corollary} on separate counters, together with the standard packages those lines need.  The retained environments print in this document as Corollary 1 in order of appearance, whereas the article prints the same items with the numbering of its own full document.  The complete native source of the article as distributed in the arXiv e-print for this exact version is bundled unchanged as \texttt{sources/169121/source0.tex}.
\begin{corollary}
\label{cor:classification}
Let $G \rightrightarrows M$ be an $s$-connected and $s$-simply connected Lie groupoid with Lie algebroid $AG$. Then the classification of "right-invariant Poisson–Nijenhuis structures" on $G$ is equivalent to the classification of "$(\Lambda, n)$-structures" on $AG$ satisfying:
\begin{itemize}
    \item $[\Lambda, \Lambda]_{\SN} = 0$,
    \item $[n, n] = 0$,
    \item $n \circ \Lambda^{\sharp} = \Lambda^{\sharp} \circ n^{*}$,
\end{itemize}
where $\Lambda \in \Gamma(\wedge^2 AG)$ and $n: \Gamma{(AG)} \to \Gamma{(AG)}$ is a vector bundle morphism covering $n_M: TM \to TM$.

Moreover, the correspondence is given explicitly by:
\[
\Pi = \overrightarrow{\Lambda}, \qquad N = \overrightarrow{n}, \qquad N_M = n_M.
\]

In other words, for a fixed Lie groupoid $G$ satisfying the above assumptions, there is a one-to-one correspondence:
\[
\left\{
\begin{array}{c}
\text{Right-invariant Poisson--Nijenhuis} \\
\text{structures on } G
\end{array}
\right\}
\quad \longleftrightarrow \quad
\left\{
\begin{array}{c}
(\Lambda, n)\text{-structures on } AG \\
\text{satisfying the three conditions}
\end{array}
\right\}.
\]
Thus, the classification problem on the global object $G$ reduces to a purely infinitesimal (often algebraic) classification problem on its Lie algebroid $AG$.
\end{corollary}

\begin{proof}
This follows directly from Theorem 13. The forward direction (from $G$ to $AG$) is given by Proposition 12, and the backward direction (from $AG$ to $G$) follows from the integration theorem for Poisson groupoids and multiplicative tensors, using the $s$-connected and $s$-simply connected assumptions. The bijectivity of the correspondence ensures that distinct $(\Lambda, n)$-structures lift to distinct PN-structures on $G$, and vice versa.
\end{proof}

\end{document}
